\section{Telescopic Series}

A series written in the form
\[
  \sum (a_{k} - a_{k+1})
\]
is called \emph{telescopic}
\mymargin{telescopic series}%
\index{series!telescopic}
because the individual terms of the sum (like the tubes of a telescope),
cancel each other out (allowing the telescope to close):
\[
  S_n 
  = \sum_{k=0}^{n-1} (a_{k} - a_{k+1})
  = \sum_{k=0}^{n-1} a_k - \sum_{k=1}^{n} a_k
  = a_0 - a_n.
\]

In theory, any series can be written in telescopic form,
given $S_n$, it suffices to choose $a_n=-S_n$, 
so that the previous relation holds. 
Writing a series in telescopic form is therefore equivalent to 
determining the sequence of partial sums.

\begin{example}[Mengoli's series]
\mymark{**}
We have
\[
  \sum_{n=1}^{+\infty} \frac{1}{n(n+1)} = 1.
\]
\end{example}
%
\begin{proof}
\mymark{**}
Indeed,
\[
  \sum_{k=1}^n \frac{1}{k(k+1)}
  = \sum_{k=1}^n \enclose{\frac{1}{k} - \frac{1}{k+1}}
  = \sum_{k=1}^n \frac{1}{k} - \sum_{k=2}^{n+1} \frac{1}{k}
  = 1 - \frac{1}{n+1} \to 1.
\]
\end{proof}